\documentclass{article}
\usepackage{graphicx} % Required for inserting images
\usepackage{booktabs}
\usepackage{amsmath}
\usepackage{amssymb}
\setlength\parindent{0pt}
\usepackage[T1]{fontenc}
% \usepackage{polski}
\title{Równania różniczkowe cząstkowe, część A}
\author{Wojciech Pachowiak}
\date{2024SZ}

\begin{document}

% \maketitle

\break

\section*{Zadanie 14}
$$ 
(e^z - e^y)\frac{\partial{u}}{\partial{x}} +
(e^x - e^z)\frac{\partial{u}}{\partial{y}} +
(e^y - e^x)\frac{\partial{u}}{\partial{z}}
= 0
$$


Układ równań charakterystyk:
\begin{align*}
    x' &= \frac{dx}{dt} = e^z - e^y  \\
    y' &= \frac{dy}{dt} = e^x - e^z  \\
    z' &= \frac{dz}{dt} = e^y - e^x
\end{align*}


Postać symetryczna:
$$
\frac{dx}{e^z - e^y} = \frac{dy}{e^x - e^z} = \frac{dz}{e^y - e^x} = dt
$$


Współczynniki nieoznaczone:
\begin{table}[h]
\begin{tabular}{@{}lllll@{}}
\toprule
                          & 0 & 0                      \\ \midrule
$\lambda_1$ & 1 & $e^x$    \\
$\lambda_2$ & 1 & $e^y$     \\
$\lambda_3$ & 1 & $e^z$    \\ \bottomrule
\end{tabular}
\end{table}

Pierwsza całka pierwsza:
$$dx + dy + dz = 0$$
$$ x+y+z = C_1 = u_1(x,y,z)  $$

Druga całka pierwsza:
$$ e^xdx + e^ydy + e^zdz = 0 $$
$$ e^x + e^y + e^z =  C_2 =  u_2(x,y,z) $$



Niezależność funkcyjna: 
$$
J =
\begin{bmatrix}
   1 & 1 & 1 \\ e^x & e^y & e^z 
\end{bmatrix}
$$

Dla $x \ne y \ne z \ne 0$ mamy $rz(J) = 2$.

\bigskip


Rozwiązanie ogólne:

$$ u(x,y,z) = F(x+y+z, e^x + e^y + e^z) $$

dla $ F \in C^1(\Omega) $.



% \input{zad16.tex}

\section*{Zadanie 17}
$$ 
x\frac{\partial{z}}{\partial{x}} +
y\frac{\partial{z}}{\partial{y}} =
z^2y
$$

$$ 
\Gamma: \quad x=t \quad y=t^2 \quad z=1 
$$

Linearyzuję:

$$ z = Z(x,y) $$
$$ u(x,y,z) = z - Z(x,y) $$

\begin{align*}
\frac{\partial{u}}{\partial{x}} &= -\frac{\partial{Z}}{\partial{x}} =
-\frac{\partial{z}}{\partial{x}} 
\\
\frac{\partial{u}}{\partial{y}} &= -\frac{\partial{Z}}{\partial{y}} =
-\frac{\partial{z}}{\partial{y}} 
\\
\frac{\partial{u}}{\partial{z}} &= 1
\end{align*}

Zlinearyzowane równanie:
$$ 
-x\frac{\partial{u}}{\partial{x}} 
-y\frac{\partial{u}}{\partial{y}} =
z^2y\frac{\partial{u}}{\partial{z}}
$$
$$ 
x\frac{\partial{u}}{\partial{x}} +
y\frac{\partial{u}}{\partial{y}} +
z^2y\frac{\partial{u}}{\partial{z}} = 0
$$


Układ równań charakterystyk:
\begin{align*}
    x' &= \frac{dx}{dt} = x  \\
    y' &= \frac{dy}{dt} = y  \\
    z' &= \frac{dz}{dt} = z^2y
\end{align*}


Postać symetryczna:
$$
\frac{dx}{x} = \frac{dy}{y} = \frac{dz}{z^2y} = dt
$$


Pierwsze z drugim:
$$
\int \frac{1}{x} dx = \int \frac{1}{y} dy 
$$
$$
\ln|x| = \ln|y| + C_1
$$
$$
\ln\left(\frac{x}{y}\right) = C_1 
$$
$$
\frac{x}{y} = C_1 = u_1(x,y,z)
$$

Drugie z trzecim:
$$
y dy = \frac{1}{z^2y} dz
$$
$$
\int dy = \int \frac{1}{z^2} dz
$$
$$
y = -\frac{1}{z} + C_2
$$
$$
y + \frac{1}{z} = C_2 = u_2(x,y,z)
$$


Niezależnośc funkcyjna: \\ 
\begin{align*}
J &= \begin{bmatrix}
    \dfrac{\partial{u_1}}{\partial{x}} & \dfrac{\partial{u_1}}{\partial{y}} & \dfrac{\partial{u_1}}{\partial{z}} \\
    \dfrac{\partial{u_2}}{\partial{x}} & \dfrac{\partial{u_2}}{\partial{y}} & \dfrac{\partial{u_2}}{\partial{z}}
\end{bmatrix} \\
&= \begin{bmatrix}
    \dfrac{1}{y} & -\dfrac{1}{y^2} & 0 \\
    0 & 1 & -\dfrac{1}{z^2}
\end{bmatrix}
\end{align*}

Całki pierwsze $u_1(x,y,z)$ i $u_2(x,y,z)$ są niezależne funkcyjnie, bo wyznacznik
$$
\begin{vmatrix}
\dfrac{1}{y} & -\dfrac{1}{y^2} \\
0 & 1
\end{vmatrix} \ne 0
$$
jest różny od zera dla dowolnego $y \in \mathbb{R}$. Zatem $rz(J) = 2$. 

\bigskip

Podstawienie na krzywej $\Gamma$:
$$
\tilde{C}_1 = \frac{t}{t^2} = \frac{1}{t}
$$
$$
t = \frac{1}{\tilde{C}_1}
$$

$$
\tilde{C}_2 = t^2 + 1
$$
$$
t = \pm \sqrt{\tilde{C}_2 - 1}
$$

$$
\Downarrow
$$

$$
\frac{1}{\tilde{C}_1} = \pm \sqrt{\tilde{C}_2 - 1}
$$
$$
\left(\frac{1}{\tilde{C}_1}\right)^2 = \tilde{C}_2 - 1
$$


Wracamy na całą dziedzinę:
$$
\left(\frac{1}{{C}_1}\right)^2 = {C}_2 - 1
$$
$$
\frac{y^2}{x^2} = y + \frac{1}{z} - 1 
$$


Szukamy postaci funkcji $z$:
$$
y + \frac{1}{z} - 1 = \frac{y^2}{x^2}
$$
$$
\frac{1}{z} = \frac{y^2}{x^2} -y + 1
$$
$$
z = \frac{1}{\dfrac{y^2}{x^2} -y + 1}
$$
$$
z = \frac{x^2}{y^2 - x^2y + x^2}
$$


Funkcja $z$ jest zdefiniowana na $\mathbb{R}^2 \setminus \{(x,y) \in \mathbb{R}^2 : y^2 - x^2y + x^2 = 0\}$.






\end{document}
